\documentclass[12pt]{article}

\usepackage[a4paper, margin=1in]{geometry}
\usepackage{amsmath}
\usepackage{graphicx}
\usepackage{booktabs}
\usepackage{hyperref}
\usepackage{float}

\title{\textbf{AI \& ML in Healthcare: NLP Research Analysis System}}
\author{Sahil Chand \and Aradhya Tiwari \and Aaryan Krishna \and Vivek Kumar Raj}
\date{}

\begin{document}

\maketitle
\tableofcontents
\newpage

% ------------------------------------------------
\section{Introduction \& Problem Statement}

The healthcare research ecosystem has witnessed exponential growth in publications related to Artificial Intelligence (AI) and Machine Learning (ML). Researchers are increasingly overwhelmed by the volume of PDF-based academic papers spanning genomics, clinical decision systems, medical imaging, and predictive healthcare analytics.

Manual analysis of such documents is:
\begin{itemize}
    \item Time-consuming
    \item Difficult to scale
    \item Prone to human bias
\end{itemize}

This project proposes a \textbf{Classical NLP-based Research Analysis System} that:
\begin{enumerate}
    \item Extracts structured information from unstructured PDF documents
    \item Removes scientific metadata noise
    \item Discovers hidden research themes using probabilistic topic modeling
    \item Measures document similarity
    \item Generates concise extractive summaries
\end{enumerate}

The system is intentionally built without Large Language Models (LLMs) to ensure interpretability and mathematical transparency.

% ------------------------------------------------
\section{Data Description}

\subsection{Dataset Source}

A curated corpus of 10 peer-reviewed research papers focusing on:
\begin{itemize}
    \item AI in Genomics
    \item Clinical Risk Prediction
    \item Machine Learning in Radiology
    \item Sensor-based Diagnostic Systems
\end{itemize}

All documents were provided in PDF format.

\subsection{Dataset Characteristics}

\begin{table}[H]
\centering
\begin{tabular}{ll}
\toprule
\textbf{Property} & \textbf{Description} \\
\midrule
Format & Multi-column PDF \\
Noise & DOIs, URLs, Emails, Citations \\
Text Artifacts & Hyphenation breaks, OCR artifacts \\
Sections Used & Abstract and Introduction \\
\bottomrule
\end{tabular}
\caption{Dataset Characteristics}
\end{table}

% ------------------------------------------------
\section{System Architecture}

The system follows a modular NLP pipeline:

\begin{verbatim}
PDF Upload
   ↓
Text Extraction (PyPDF2)
   ↓
Regex Cleaning
   ↓
Tokenization & Lemmatization (NLTK)
   ↓
TF-IDF Vectorization
   ↓
LDA Topic Modeling
   ↓
Cosine Similarity Matrix
   ↓
Extractive Summarization
   ↓
Streamlit UI
\end{verbatim}

% ------------------------------------------------
\section{Methodology}

\subsection{Text Preprocessing}

The preprocessing pipeline consists of:
\begin{itemize}
    \item Lowercasing
    \item Removal of URLs, DOIs, citation indices
    \item Stopword filtering
    \item Lemmatization using WordNet
    \item Removal of short tokens ($<$ 3 characters)
\end{itemize}

% ------------------------------------------------
\subsection{Feature Engineering: TF-IDF}

We transform text into numerical vectors using Term Frequency–Inverse Document Frequency:

\begin{equation}
TFIDF(t, d) = TF(t, d) \times \log\left(\frac{N}{DF(t)}\right)
\end{equation}

Where:
\begin{itemize}
    \item $TF(t, d)$ = frequency of term $t$ in document $d$
    \item $DF(t)$ = number of documents containing term $t$
    \item $N$ = total number of documents
\end{itemize}

Configuration:
\begin{itemize}
    \item min\_df = 2
    \item max\_df = 0.85
    \item ngram\_range = (1,2)
\end{itemize}

% ------------------------------------------------
\subsection{Topic Modeling: Latent Dirichlet Allocation}

LDA assumes each document is a mixture of topics and each topic is a distribution over words.

\begin{equation}
P(w|d) = \sum_{k=1}^{K} P(w|z_k) \cdot P(z_k|d)
\end{equation}

Where:
\begin{itemize}
    \item $z_k$ represents topic $k$
    \item $P(w|z_k)$ is the word distribution per topic
    \item $P(z_k|d)$ is the topic distribution per document
\end{itemize}

We selected $K = 4$ topics for optimal interpretability.

% ------------------------------------------------
\subsection{Cosine Similarity}

To measure document similarity:

\begin{equation}
\cos(\theta) = \frac{A \cdot B}{||A|| \cdot ||B||}
\end{equation}

Values range from 0 (completely different) to 1 (identical).

% ------------------------------------------------
\subsection{Extractive Summarization}

The summarization process:
\begin{enumerate}
    \item Extracts the Abstract section
    \item Tokenizes sentences
    \item Ranks sentences using TF-IDF weights
    \item Selects top 3 highest scoring sentences
    \item Reorders them chronologically
\end{enumerate}

This ensures logical coherence without hallucination.

% ------------------------------------------------
\section{Evaluation \& Results}

\subsection{Identified Topics}

\begin{table}[H]
\centering
\begin{tabular}{ll}
\toprule
\textbf{Topic} & \textbf{Identified Theme} \\
\midrule
Topic 1 & Genomics \& Gene Analysis \\
Topic 2 & Clinical Risk Prediction \\
Topic 3 & Algorithm Development \& Ethics \\
Topic 4 & Sensor-based Diagnostics \\
\bottomrule
\end{tabular}
\caption{LDA Identified Topics}
\end{table}

The topics displayed strong semantic cohesion across documents.

% ------------------------------------------------
\section{Optimization Steps}

\begin{itemize}
    \item Replaced spaCy with NLTK to reduce memory footprint
    \item Added regex-based hyphen repair
    \item Expanded domain stopword list
    \item Pinned dependency versions for cloud stability
\end{itemize}

% ------------------------------------------------
\section{Limitations}

\begin{itemize}
    \item Bag-of-words assumption ignores word order
    \item No semantic embeddings
    \item Limited cross-document reasoning
    \item No real-time web retrieval
\end{itemize}

% ------------------------------------------------
\section{Conclusion}

This project demonstrates how classical NLP techniques can structure and analyze healthcare research documents. By combining TF-IDF, LDA, cosine similarity, and extractive summarization, we transformed unstructured PDFs into interpretable thematic insights.

The system fulfills all Milestone 1 objectives while maintaining transparency and computational efficiency.

% ------------------------------------------------
\section*{References}

\begin{enumerate}
    \item Blei, D., Ng, A., \& Jordan, M. (2003). Latent Dirichlet Allocation. Journal of Machine Learning Research.
    \item Manning, C., Raghavan, P., \& Schütze, H. (2008). Introduction to Information Retrieval.
    \item Pedregosa et al. (2011). Scikit-learn: Machine Learning in Python.
    \item Bird, S., Klein, E., \& Loper, E. (2009). Natural Language Processing with Python.
    \item Salton, G., \& McGill, M. (1983). Introduction to Modern Information Retrieval.
\end{enumerate}

\end{document}
